\documentclass[conference]{IEEEtran}
\usepackage[T1]{fontenc}
\usepackage[utf8]{inputenc}
\usepackage{amsmath,amssymb,booktabs}
\usepackage[hidelinks]{hyperref}
\title{Sparse-Channel EEG Affect Representation Under Cross-Dataset and Montage Shift}
\author{\IEEEauthorblockN{Author information to be confirmed}\IEEEauthorblockA{Institutional affiliation to be confirmed}}
\begin{document}
\maketitle
\begin{abstract}
Portable affective EEG decoding is challenged by sensor density, electrode reference, participant identity and acquisition-domain shift. This work specifies a controlled benchmark that separates the effect of channel sparsity from corpus and montage incompatibility. The proposed design uses subject-disjoint evaluation, auditable channel coordinates, train-only preprocessing, strong low-capacity baselines and uncertainty at participant level. Cross-corpus inference is permitted only when acquisition and label compatibility are established. This document is a methods-stage manuscript draft: no dataset rights, experiments or predictive results are claimed as completed.
\end{abstract}
\begin{IEEEkeywords}
EEG, domain shift, montage, calibration, participant-disjoint evaluation
\end{IEEEkeywords}
\section{Introduction}
Low-electrode-count EEG systems reduce acquisition and deployment burden, yet they do not necessarily preserve information acquired by research-grade montages. Apparent emotion-decoding accuracy can arise from stable participant or recording signatures rather than transferable affect-related information. A recent analysis of DEAP and DREAMER found substantial performance reductions after enforcing participant independence \cite{pandilova2026}. These results motivate evaluation at the participant level rather than treating overlapping EEG windows as independent observations.

Prior work has addressed transfer between subjects, sessions and datasets \cite{review2024,crossadapt2025}. More specifically, consumer-grade, cross-dataset EEG studies already compare feature engineering and neural methods \cite{consumer2025}. Consequently, sparse-channel decoding or EEG domain adaptation alone cannot constitute a credible novelty claim.

\textbf{Research question:} how much predictive utility and calibration are lost when an EEG model is transferred across participants and acquisition domains while reducing the available electrode set, after separating changes in channel count from changes in electrode position and reference?

The proposed differentiator is a factorial \emph{measurement compatibility} audit: (i) legal access and physical electrode correspondence, (ii) reference and label compatibility, (iii) subject-disjoint versus domain-held-out evaluation and (iv) calibration, including failures and negative controls. The intent is to quantify limits and failure modes, not to claim equivalence with a particular consumer headset.

\textbf{Scope and contribution status.} This is a planned benchmark protocol and a manuscript draft. Its technical contribution and novelty must be evaluated against the cited studies with actual reproduced baselines before any publication claim. No new human participants are needed; the candidate datasets contain existing recordings subject to their own reuse conditions.

\section{Materials and Methods}
\subsection{Eligibility, rights and metadata}
DEAP \cite{deap2012}, DREAMER \cite{dreamer2018} and AMIGOS \cite{amigos2021} are candidate datasets, not presently approved inputs. As a rights-first alternative, GAMEEMO v3 \cite{gameemo2020} (public Mendeley Data, CC BY 4.0, 14-channel EPOC+ with SAM ratings) and FACED NEMAR v1.1.4 \cite{faced2023} (public BIDS derivative, CC BY 4.0, 32-channel EEG with rated stimuli) are also considered, with versions and attribution stated in the dataset registry. Neither alternative has been downloaded or validated. GAMEEMO gaming tasks and FACED video tasks may differ substantially in label semantics, recording context and population. Before analysis, record the dataset version, approval or licence, release files, cryptographic checksums, unique participants, sessions/trials, sampling frequency, EEG channel names, electrode coordinates, reference scheme, filter history and target-label definition. Exclude any dataset lacking legally usable original signals or sufficient montage/label metadata. Do not substitute unofficial mirrors. Existing participant data must not be committed to source control.

\subsection{Montage compatibility}
For each acquisition system define the channel set $C_d$, electrode positions $\mathbf{r}_{d,c}$ and reference operator $R_d$. A candidate sparse montage $S_d \subseteq C_d$ is admissible only if its electrode positions and reference can be documented; the target count $k\in\{2,4,8\}$ is conditional on actual electrodes. The full acquisition set is the upper-information comparator. Identically named electrodes are not treated as equivalent after an unverified reference transformation. If two datasets lack defensible common sensors/references, cross-corpus transfer for that pair is declared \emph{not estimable}; within-dataset results may still be reported. Neither a DEAP four-channel subset nor a DREAMER subset is labelled a Muse S Athena measurement.

\subsection{Targets, preprocessing and splitting}
Primary target is a valence dimension only if its original scale and timing are documented; arousal is secondary. Original continuous ratings are preserved. Any binary threshold must be defined \emph{a priori} from scale semantics or from training data only, with exact threshold recorded; test-label medians are prohibited. The unit of independence is the participant; epochs from one trial must not cross a fold boundary.

Outer evaluation is leave-one-participant-out within each eligible corpus. Training subjects alone determine resampling, referencing, bandpass design, artifact thresholds, feature selection, missing-value imputation, scaling, calibration and hyperparameter search, with group-aware inner folds. Deterministic seeds and exact split manifests are saved. For cross-corpus analysis, the target dataset remains entirely unseen while model selection is performed on source subjects; a target-domain adaptation variant is an explicitly separate transductive analysis, never mixed with inductive transfer.

\subsection{Models and controls}
Comparators include a majority/constant predictor, ridge or elastic-net logistic regression on log-bandpower, a compact convolutional model when justified by sample size and a full-channel reference. The sparse-channel proposal must be compared with a \emph{capacity-matched} sparse baseline using the same train folds and preprocessing. Negative controls are a within-training-fold label permutation and a non-EEG participant-prior diagnostic where applicable. Shared train/test participants in a deliberately leaky diagnostic may demonstrate bias but must never enter the primary result.

\subsection{Endpoints and inference}
For categorical targets, preregister participant-aggregated balanced accuracy as the primary performance metric and Brier score as the primary calibration diagnostic; AUROC, macro-F1 and ECE are secondary with reported prevalence. Let $\Delta_{d,k}$ be paired full-minus-sparse balanced accuracy, computed per held-out participant. Compare $\Delta_{d,k}$ across eligible $k$ and domains, without selecting the best value on the test set. For a prespecified paired comparison, calculate one Brier score per held-out participant and model after aggregating window probabilities within each annotated trial. Resample whole participants with replacement while keeping each participant's reference and candidate Brier scores paired. Report the equally participant-weighted difference (reference minus candidate; positive means lower candidate Brier) with a 95\% percentile bootstrap interval, its seed, number of independent participants and number of resamples. This interval cannot make a small or dependent cohort representative; if sample or model constraints invalidate the interval, report raw per-participant effects and an alternative defensible uncertainty analysis. Label-conditioned experiments and multiple target/dataset comparisons are identified as exploratory unless a primary family is frozen before test evaluation.

\subsection{Reproducibility and stop rules}
Release preprocessing configurations, source hashes, split assignments that reveal no private identity, feature and model versions, machine/software metadata and figure scripts. Stop cross-corpus claims if montage/reference compatibility or annotation semantics fail. Stop any accuracy claims if complete participant separation or baseline reproduction cannot be demonstrated. No real Muse or Tobii raw streams are part of this study under current access conditions.

\section{Results: Evidence Required}
\textbf{No empirical results are available in this manuscript version.} This section is an explicit reporting contract, not a table of observations.

\begin{table}[htbp]
\caption{Primary evaluation table, to be populated only after the data and evidence gates pass}
\centering
\begin{tabular}{llll}
\toprule
Corpus & Montage & Balanced acc. & Brier \\
\midrule
Eligible corpus & Full & Pending & Pending\\
Eligible corpus & Sparse $k=4$ & Pending & Pending\\
Held-out corpus & Compatible subset & Not yet eligible & Pending\\
\bottomrule
\end{tabular}
\end{table}

The final report will contain participant-level paired differences and intervals, the exact number of legal datasets and unique participants, montage/reference exclusions, train-only preprocessing audit, loss of utility versus electrode count and a corpus-by-montage comparison only where scientifically identifiable. Negative controls and failure cases must be shown alongside positive results. Any claim unsupported by an actual data and code run will be omitted.

\section{Discussion and Limitations}
The design intentionally separates the measurement problem from the modelling problem. A low-channel representation can appear robust because preprocessing inadvertently uses target-subject statistics, because participants leak across windows, or because different corpus labels are mapped to superficially similar categories. Each failure invalidates a broad generalization claim. Recent evidence of severe identity confounding in qEEG benchmarks makes these checks indispensable \cite{pandilova2026}.

The benchmark cannot establish physiological equivalence between different hardware, and affective ratings are operational reference annotations rather than direct neural ground truth. Potentially small eligible corpus overlap may prevent a cross-dataset factorial test; that is a scientifically interpretable feasibility outcome, not a reason to interpolate incompatible sensors. A negative or near-chance external result is publishable only insofar as the protocol is sufficiently rigorous, novel and transparent.

\section{Conclusion}
The proposed study is intended to measure the limits of sparse EEG transfer under explicitly audited montage, person and domain changes. Completion depends on legitimate dataset access, frozen analysis decisions and reproducible baselines. No clinical or diagnostic inference is proposed.

\bibliographystyle{IEEEtran}
\bibliography{references}
\end{document}
